\clearpage
\begingroup
\tcbset{myboxstyle/.append style={breakable=false}}
\everymath{\textstyle}
\subsection{Математическое ожидание функции случайных величин}

Пусть $f\colon\R^2\to\R$ --- борелевская функция.

\begin{theorem}[Дискретное совместное распределение]
    Если $p_{ij}=\Pr\{\xi=x_i,\eta=y_j\}$ и
    $\sum_{i,j}|f(x_i,y_j)|p_{ij}<\infty$, то
    \[
        \EE f(\xi,\eta)=\sum_{i,j}f(x_i,y_j)p_{ij}.
    \]
\end{theorem}

\begin{theorem}[Абсолютно непрерывное совместное распределение]
    Если $(\xi,\eta)$ имеет плотность $p_{\xi,\eta}$ и
    \[
        \iint_{\R^2}|f(x,y)|p_{\xi,\eta}(x,y)\,\diff x\,\diff y<\infty,
    \]
    то
    \[
        \EE f(\xi,\eta)
        =\iint_{\R^2}f(x,y)p_{\xi,\eta}(x,y)\,\diff x\,\diff y.
    \]
\end{theorem}

\subsection{Ковариация}

\begin{definition}
    Для случайных величин $\xi$ и $\eta$ с конечными вторыми моментами
    \textbf{ковариацией} называется число
    \[
        \operatorname{cov}(\xi,\eta)
        =\EE\bigl[(\xi-\EE\xi)(\eta-\EE\eta)\bigr].
    \]
\end{definition}

\begin{theorem}[Свойства ковариации]
    При конечности вторых моментов справедливы следующие свойства.
    \begin{enumerate}
        \item $\operatorname{cov}(\xi,\xi)=\DD\xi$.
        % Дополнение из лекции 17, с. 5.
        \item $\operatorname{cov}(\xi,\eta)=\EE(\xi\eta)-\EE\xi\cdot\EE\eta$.
        \item Ковариация является билинейной формой: для $a,b\in\R$
        \[
            \begin{aligned}
                \operatorname{cov}(a\xi+b\eta,\zeta)
                &=a\operatorname{cov}(\xi,\zeta)+b\operatorname{cov}(\eta,\zeta),\\
                \operatorname{cov}(\zeta,a\xi+b\eta)
                &=a\operatorname{cov}(\zeta,\xi)+b\operatorname{cov}(\zeta,\eta).
            \end{aligned}
        \]
        \item Из независимости $\xi$ и $\eta$ следует
        $\operatorname{cov}(\xi,\eta)=0$. Обратное утверждение неверно.
    \end{enumerate}
\end{theorem}
% Дополнение из лекции 17, с. 5.
\begin{theorem}[Дисперсия суммы]
    Если $\xi$ и $\eta$ имеют конечные вторые моменты, то
    \[
        \DD(\xi+\eta)=\DD\xi+2\operatorname{cov}(\xi,\eta)+\DD\eta.
    \]
    В частности, для независимых $\xi$ и $\eta$
    \[
        \DD(\xi+\eta)=\DD\xi+\DD\eta.
    \]
\end{theorem}

\begin{Proof}
    По линейности математического ожидания
    \[
        \begin{aligned}
            \DD(\xi+\eta)
            &=\EE\bigl[(\xi-\EE\xi)+(\eta-\EE\eta)\bigr]^2\\
            &=\EE(\xi-\EE\xi)^2
              +2\EE\bigl[(\xi-\EE\xi)(\eta-\EE\eta)\bigr]
              +\EE(\eta-\EE\eta)^2\\
            &=\DD\xi+2\operatorname{cov}(\xi,\eta)+\DD\eta.
        \end{aligned}
    \]
    При независимости ковариация равна нулю.
\end{Proof}
\endgroup
